%% =====================================================================
%%  2 / 7   Literature Review  文献综述                         成员 2
%% =====================================================================
\section{Literature Review}
\label{sec:litreview}

\subsection{From clinical decision support to medical agents}

For most of the past decade, clinical AI followed an advisory division of labor. \citet{topol2019} described high-performance medicine as a partnership in which deep learning handles pattern recognition and data sorting while physicians make diagnostic and treatment decisions. Early discussions of clinical autonomy assumed the same baseline. \citet{bitterman2020autonomy} argued that medical algorithms should first prove safe as assistive tools and demonstrate benefit on prospective clinical outcomes across diverse patient groups, and warned that universal full autonomy across clinical settings was unlikely in the near term.

LLM-based medical agents depart from that advisory setup by linking multi-turn planning with persistent memory and external tool calls across an episode of care \citep{wang2025baymax}. In \citet{li2024agenthospital}, LLM doctor, nurse, and patient agents interact across a simulated hospital so clinical policies can be refined offline without putting real patients at risk. Other engineering work automates workflow construction itself: \citet{zhuang2025learning} represent clinical multi-agent systems as searchable graphs that combine conditional branching with self-reflection.

Once benchmarks moved from static licensing exams into interactive clinical environments, performance limits became far clearer. \citet{schmidgall2024agentclinic} evaluated multimodal LLM agents in sequential simulated consultations and found that accuracy drops steeply when models must gather patient histories over multiple turns. Injecting realistic cognitive biases into either the patient or doctor agent lowered diagnostic accuracy further and reduced simulated patient compliance. Inside a virtual electronic health record, \citet{jiang2025medagentbench} tested whether agents can both retrieve patient data and execute write operations across 300 physician-authored tasks; even the highest-scoring model failed roughly three tasks in ten. Benchmarks in this literature judge whether an agent finishes a clinical task in simulation. They leave unexamined when a clinician should delegate execution to a system with a substantial failure rate, or how a reviewer can check a multi-step action sequence before harm occurs.

\subsection{Human oversight, automation bias, and skill erosion}

Empirical studies of clinicians using AI show that keeping a human in the loop does not reliably catch machine errors or improve clinical judgement. Even with advisory models, collaboration often falls short of standalone model scores. In a single-blind randomized clinical trial of 50 internal, family, and emergency medicine physicians, \citet{goh2024diagnostic} compared diagnostic reasoning on complex case vignettes with and without access to GPT-4 alongside standard clinical references. Physicians given LLM access achieved a median score of 76\%, almost indistinguishable from the 74\% scored with conventional resources alone ($p = 0.60$), even though GPT-4 scored higher on its own in exploratory analyses.

Human-factors research identifies structural reasons why nominal supervision breaks down under clinical workload. \citet{bainbridge1983} observed that automating routine work often leaves people with infrequent, harder monitoring and exception handling when automation fails. High reliability compounds this effect: \citet{parasuraman2010} showed that dependable automation induces complacency and narrows attention, so operators miss rare errors more often. Earlier human-factors work treated automation bias mainly as a byproduct of multitasking, but a systematic review of 40 studies by \citet{lyell2017automation} found that it also occurs in single diagnostic tasks whenever verification complexity and cognitive load are high. That mechanism bears directly on medical agents, since checking a trail of tool calls and draft orders demands far more effort than verifying one diagnostic suggestion.

Routine reliance on AI can also degrade unaided clinical skill. In a multicenter observational study of colonoscopy, \citet{budzyn2025deskilling} found that endoscopists' unaided adenoma detection rate fell from 28.4\% (226 of 795 procedures) before regular AI exposure to 22.4\% (145 of 648 procedures, $p = 0.0089$) afterward. In a scoping review across medical specialties, \citet{natali2025deskilling} organized the mechanisms through which habitual AI use weakens diagnostic and procedural competence, while \citet{hunter2024reliance} proposed periodic reliance drills to test whether clinicians still catch errors when automated support is removed. Patient attitudes reflect similar sensitivity around human control: in a multinational hospital survey, \citet{busch2025attitudes} found that acceptance of medical AI varies substantially with the clinical task, the clinical stakes, and the extent of physician involvement. Almost all of this empirical evidence, however, comes from diagnostic advice or image detection. Little is known about how clinical oversight and physician skill hold up once software writes directly to patient records.

\subsection{Autonomy taxonomies, liability, and clinical governance}

Graded autonomy frameworks attempt to replace coarse distinctions between human and automated care. \citet{ye2025autonomy} built on earlier medical autonomy models \citep{bitterman2020autonomy} to divide medical LLMs into four tiers, from informational tools (L0) and decision support through supervised autonomous agents (L3), and matched evaluation expectations to each tier. Classifying capability and risk clarifies what a model can attempt, yet it does not specify the clinical conditions under which a physician or hospital may authorize a particular action.

Legal and safety analyses point to institutional barriers that sit behind those technical tiers. In an analysis of United States medical malpractice law, \citet{price2019liability} showed that tort doctrine shields physicians who follow the prevailing standard of care. That rule gives risk-averse clinicians an incentive to treat AI mainly as a confirmatory check and reject unfamiliar recommendations, even when following the system might benefit a patient. From a safety-engineering perspective, \citet{habli2020accountability} contended that static pre-market approval cannot sustain accountability for adaptive clinical AI; because responsibility spans clinicians, developers, and safety engineers, safety assurance must continue after deployment as systems encounter shifting patient populations.

Recent clinical guidance focuses on professional roles and workplace training. \citet{clayton2025clinicians} held that even as software absorbs data sorting and pattern synthesis, diagnostic reasoning and patient communication remain core physician responsibilities. \citet{abdulnour2025education} translated that position into supervision curricula designed to teach clinicians how to verify AI outputs critically and guard against deskilling, mis-skilling, and never-skilling. \citet{ong2026flight} drew on commercial aviation to propose explicit ``flight rules'' for clinical AI, pairing rehearsed collaboration procedures with incident reporting and minimum unaided practice requirements so clinicians stay in active command. Together, these proposals target system-level regulation or long-term clinical habits, leaving undefined what a clinician must specify before authorizing a single delegated action.

\subsection{Research gap}

Research on medical AI remains divided across disciplinary lines. Engineering work advances multi-turn planning and tool execution in simulated clinics and virtual health records, while still judging systems primarily by benchmark completion. Human-factors and clinical studies document automation bias and skill erosion in real and simulated care, yet their evidence base centers on advisory tools in which clinicians still carry out the final act. Governance frameworks, malpractice analyses, and aviation-inspired protocols map autonomy tiers and professional duties at the level of systems or training programs, without specifying what a clinician must establish before delegating a concrete clinical action to an agent. What remains missing is an account of the authorization relationship itself once medical AI moves from recommending to acting: the conditions under which clinical execution can be legitimately delegated, and how crossing from advice to action alters clinician and patient judgements of trust, oversight, and accountability.
